\documentclass[10pt]{article}
\usepackage[T1]{fontenc}
\usepackage[utf8]{inputenc}
\usepackage[french]{babel}
\usepackage{amsmath,amssymb,mathrsfs,mathpazo}
\usepackage{geometry}
\geometry{a4paper,margin=2cm}
\usepackage{graphicx,array,tabularx,tabularray,colortbl,multirow,multicol,booktabs,enumitem,pifont,ulem,fancybox,diagbox,cancel}
\usepackage{pstricks,pstricks-add,pst-plot,pst-tree,pst-node,pst-text,pst-func,pst-math,pst-bspline,pst-3dplot,pst-eucl,pst-solides3d}
\let\pscalculate\undefined
\usepackage{pst-bezier}
\usepackage{tikz,pgfplots,tkz-tab}
\usetikzlibrary{arrows,arrows.meta,patterns,calc,positioning,trees,decorations.markings,shapes.geometric,shapes.misc,shapes.arrows,intersections,decorations.pathmorphing,decorations.pathreplacing,backgrounds,shadows,babel}
\pgfplotsset{compat=1.18}
\usepgfplotslibrary{fillbetween}
\usepackage[np]{numprint}
\usepackage{listings,makecell,pas-tableur,twoopt}
\let\barre\undefined
\usepackage[pstricks]{tablvar}
\definecolor{darkgreen}{RGB}{0,100,0}
\providecommand{\red}{\color{red}}
\providecommand{\blue}{\color{blue}}
\providecommand{\green}{\color{green}}
\providecommand{\black}{\color{black}}
\providecommand{\white}{\color{white}}
\providecommand{\gray}{\color{gray}}
\providecommand{\cyan}{\color{cyan}}
\providecommand{\magenta}{\color{magenta}}
\providecommand{\yellow}{\color{yellow}}
\providecommand{\violet}{\color{violet}}
\providecommand{\orange}{\color{orange}}
\providecommand{\euro}{\texteuro}
\providecommand{\EUR}{\texteuro}
\providecommand{\ieme}{\textsuperscript{e}}
\providecommand{\degres}{\ensuremath{^\circ}}
\providecommand{\degree}{\ensuremath{^\circ}}
\providecommand{\vv}[1]{\overrightarrow{#1}}
\providecommand{\Coord}[1]{#1}
\providecommand{\decosix}{}
\providecommand{\starredbullet}{\ensuremath{\bullet}}
\pagestyle{empty}
\setlength{\parindent}{0pt}
\usepackage{adjustbox}\DeclareUnicodeCharacter{2236}{:}\DeclareUnicodeCharacter{2190}{\ensuremath{\leftarrow}}
\providecommand{\R}{\mathbb{R}}
\providecommand{\N}{\mathbb{N}}
\providecommand{\Z}{\mathbb{Z}}
\providecommand{\Q}{\mathbb{Q}}
\providecommand{\D}{\mathbb{D}}
\providecommand{\C}{\mathbb{C}}
\providecommand{\e}{{\mathrm{e}}}
\providecommand{\vect}[1]{\overrightarrow{#1}}
\providecommand{\vectt}[1]{\overrightarrow{\mathrm{#1}}}
\providecommand{\vvt}[1]{\overrightarrow{\mathrm{#1}}}
\providecommand{\np}[2][]{{#2~#1}}
\providecommand{\vv}[1]{{\overrightarrow{#1}}}
\renewcommand{\d}{{\mathrm{d}}}
\providecommand{\degres}{{^\circ}}
\providecommand{\degre}{{^\circ}}
\providecommand{\blue}{\color{blue}}
\providecommand{\red}{\color{red}}
\providecommand{\enskip}{\quad}
\providecommand{\llbracket}{\lbrack\!\lbrack}
\providecommand{\rrbracket}{\rbrack\!\rbrack}
\providecommand{\euro}{{\text{€}}}
\providecommand{\no}{{\text{n°}}}
\providecommand{\interval}[2]{{[#1;#2]}}
\providecommand{\intervFO}[2]{{[#1;#2[}}
\providecommand{\dotfill}{\ldots}

\providecommand{\ds}{\displaystyle}
\providecommand{\textding}[1]{#1}
\providecommand{\Oij}{\(\mathrm{O};\vec{\imath},\vec{\jmath}\)}
\providecommand{\Oijk}{\(\mathrm{O};\vec{\imath},\vec{\jmath},\vec{k}\)}
\providecommand{\Ouv}{\(\mathrm{O};\vec{u},\vec{v}\)}
\providecommand{\intervFF}[2]{\ensuremath{[#1;#2]}}
\renewcommand{\euro}{\ifmmode\text{\texteuro}\else\texteuro\fi}
\DeclareUnicodeCharacter{00B0}{\ifmmode^{\circ}\else\textdegree\fi}
\DeclareUnicodeCharacter{20AC}{\ifmmode\text{\texteuro}\else\texteuro\fi}
% Énoncés : Éducation nationale. Transcription et corrigés : APMEP, Denis Vergès et contributeurs. Découpage et classement : CoopMaths, Sébastien Lozano, Éric Elter et contributeurs. Adaptation de la présentation : Maths Explorer.
% CC BY-SA : https://coopmaths.fr/www/about/
% https://coopmaths.fr/alea/static/eam/2026/tex/eam_specialite_2026_06_etrangers_3.tex
\begin{document}
\noindent\textit{APMEP / CoopMaths. CC BY-SA. Présentation adaptée par Maths Explorer.}\par\medskip
Soit $f$ une fonction définie sur $]0\ ;\ +\infty[$.

On a tracé dans le repère ci-dessous la courbe $C_f$ représentative de la fonction $f$ et la droite $T_A$ tangente à la courbe $C_f$ au point $A$ de coordonnées $(1\ ;\ 20)$.

La droite $T_A$ passe par le point $B$ de coordonnées $(3\ ;\ 10)$.


\begin{center}
	\begin{tikzpicture}[x=1.5cm,y=0.15cm]
		
		
		
		
		
		\foreach \x in {0,1,2,3,4,5,6,7}
		{
			\draw[dotted] (\x,0) -- (\x,55);
		}
		
		
		\foreach \y in {0,10,20,30,40,50}
		{
			\draw[dotted] (0,\y) -- (7.5,\y);
		}
		
		
		\draw[->,thick] (0,0) -- (7.7,0);
		\draw[->,thick] (0,0) -- (0,56);
		
		
		\foreach \x in {0,1,2,3,4,5,6,7}
		{
			\draw (\x,0) -- (\x,-1);
			\node[below] at (\x,-1.2) {\small $\x$};
		}
		
		
		\foreach \y in {10,20,30,40,50}
		{
			\draw (0,\y) -- (-0.08,\y);
			\node[left] at (-0.08,\y) {\small $\y$};
		}
		
		
		\draw[very thick, domain=0.18:7.4, samples=300, smooth, variable=\x]
		plot ({\x},{(4*\x*\x+7*\x+9)/\x});
		
		
		\draw[thick, domain=0:5.4, samples=2, variable=\x]
		plot ({\x},{-5*\x+25});
		
		
		\fill (1,20) circle (1.5pt);
		\fill (3,10) circle (1.5pt);
		
		
		\node[above right] at (1,20) {$A$};
		\node[above right] at (3,10) {$B$};
		\node at (0.45,47) {$C_f$};
		\node at (4.1,7) {$T_A$};
		
	\end{tikzpicture}
\end{center}

\begin{enumerate}
\item \begin{enumerate}
\item Donner $f(1)$.
\item Déterminer la valeur de $f'(1)$.
\item Justifier que l'équation réduite de la tangente $T_A$ à $C_f$ au point $A$ est :
\[
y=-5x+25.
\]
\end{enumerate}
\end{enumerate}

Pour la suite de l'exercice, la fonction $f$ est définie sur $]0\ ;\ +\infty[$ par $f(x)=\dfrac{4x^2+7x+9}{x}$.

On admet que $f$ est dérivable sur cet intervalle.

\begin{enumerate}
\setcounter{enumi}{1}
\item \begin{enumerate}
\item Démontrer que pour tout $x\in]0\ ;\ +\infty[$, on a $f'(x)=\dfrac{(2x-3)(2x+3)}{x^2}$.
\item Étudier le signe de $f'(x)$ sur $]0\ ;\ +\infty[$.
\item En déduire les variations de la fonction $f$ sur $]0\ ;\ +\infty[$.
\end{enumerate}
\item Existe-t-il une tangente à $C_f$ parallèle à la droite d'équation $y=3x+5$ ?

Justifier votre réponse.
\end{enumerate}

\end{document}
